\documentclass[11pt,a4paper]{article}
\usepackage[margin=25mm]{geometry}
\usepackage{amsmath,amssymb,amsthm}
\usepackage[T1]{fontenc}
\usepackage{lmodern}
\usepackage[hidelinks]{hyperref}
\setlength{\emergencystretch}{2em}
\newtheorem{lemma}{Lemma}
\newtheorem{proposition}[lemma]{Proposition}
\newcommand{\E}{\mathbb E}
\newcommand{\Pp}{\mathbb P}
\newcommand{\CE}{C_{\mathrm E}}
\title{\vspace{-1.5em}Smoothing and first-counterexample scalar gates}
\author{}
\date{27 October 2026}
\begin{document}
\maketitle
\vspace{-2.5em}
\noindent\textbf{Scope.}
This note independently checks the scalar step of a proposed proof redesign.
It does not establish the acquisition or local-exclusion theorems, verify their
certificates, or authorize removal of a proof component. Baseline:
\texttt{6346fc1b4432bc5982822aa696eb566ed9e1e1ee} in the fork
\texttt{kolodziejek/berry-esseen-certificates}.

Let $\mathcal P$ be all real laws with mean zero, variance one and finite
$\rho(P)=\E_P|X|^3$. For iid $X_i$ with law $P$, put
\[
 S_n=\sum_{i=1}^nX_i,\qquad
 \Delta_n(P)=\sup_{H\in\mathbb R}
       |\Pp_P(S_n\le H)-\Phi(H/\sqrt n)|.
\]
\[
 C_n=\sup_{P\in\mathcal P}\frac{\sqrt n\,\Delta_n(P)}{\rho(P)},\qquad
 \CE=\frac{3+\sqrt{10}}{6\sqrt{2\pi}}.
\]
Thus $H$ always denotes a raw sum threshold. All terminating decimals below
denote exact rationals. Lyapunov's inequality gives $\rho(P)\ge1$.

\begin{lemma}[Support-free smoothing]\label{lem:smooth}
For every integer $n\ge2$,
\begin{equation}\label{eq:smooth}
 C_n\le\sqrt{\frac n{n-1}}\,C_{n-1}
       +\frac{\sqrt n}{15(n-1)^{3/2}}
       +\frac{2\sqrt n}{25(n-1)^2}.
\end{equation}
No bounded support, extremizer or higher moment is assumed.
\end{lemma}
\begin{proof}
Fix $P\in\mathcal P$ and $H\in\mathbb R$, write $v=n-1\ge1$, and set
$F_v(t)=\Phi(t/\sqrt v)$. Conditioning on $X_n$ and applying the defining
bound for $C_v$ to the same marginal at every threshold gives
\begin{equation}\label{eq:condition}
 |\Pp_P(S_n\le H)-\E F_v(H-X_n)|\le C_v\rho/\sqrt v.
\end{equation}
Only the smooth Gaussian function is expanded. Taylor's theorem gives
\[
 |F_v(H-y)-F_v(H)+yF'_v(H)-\tfrac12 y^2F''_v(H)|
 \le\frac{\|\phi''\|_\infty}{6v^{3/2}}|y|^3.
\]
In fact $\|\phi''\|_\infty=\phi(0)$: since
$\phi''(z)=(z^2-1)\phi(z)$, the assertion is immediate on $z^2\le1$;
on $u=z^2\ge1$, $(u-1)e^{-u/2}$ is maximized at $u=3$, with value
$2e^{-3/2}<1$. Taking expectation therefore yields
\begin{equation}\label{eq:spatial}
 |\E F_v(H-X_n)-F_v(H)-\tfrac12F_v''(H)|
 \le\frac{\phi(0)\rho}{6v^{3/2}}<\frac{\rho}{15v^{3/2}}.
\end{equation}
The elementary bound $\phi(0)<2/5$ is established below.

Direct differentiation in the variance gives
$\partial_vF_v=F_v''/2$ and $\partial_v^2F_v=F_v^{(4)}/4$, where primes
differentiate the raw argument. Hence
\[
 F_{v+1}(H)-F_v(H)-\tfrac12F_v''(H)
 =\frac14\int_0^1(1-t)F_{v+t}^{(4)}(H)\,dt.
\]
Since $F_w^{(4)}(H)=w^{-2}\phi'''(H/\sqrt w)$, its absolute value is at most
$\|\phi'''\|_\infty/(8v^2)$. The elementary Gaussian Fourier identity,
differentiated under its absolutely convergent integral, implies
\[
 \|\phi'''\|_\infty\le
 \frac1{2\pi}\int_{\mathbb R}|t|^3e^{-t^2/2}\,dt=\frac2\pi.
\]
Indeed the positive-half-line integral is $2$, by $u=t^2/2$.
The heat remainder is consequently at most
$1/(4\pi v^2)<2/(25v^2)$. Combine this with
\eqref{eq:condition}--\eqref{eq:spatial}, multiply by $\sqrt n/\rho$,
and use $\rho\ge1$. The resulting estimate is uniform in $H$ and $P$,
so both suprema are legitimate.

For a strict event, let $H_j\uparrow H$: its probability is the limit of the
closed-event probabilities, while $\Phi$ is continuous. Thus the same
bound covers both jump edges and reflection. Neither an indicator nor
the distribution function of $S_v$ was differentiated. Rare remote atoms
need only the assumed third absolute moment.
\end{proof}

\begin{lemma}[Elementary constants]\label{lem:constants}
One has $\pi>157/50>25/8$ and $0<\CE<41/100$.
\end{lemma}
\begin{proof}
Let $a=\arctan(1/5)$ and $b=\arctan(1/239)$. Tangent addition gives
$\tan(2a)=5/12$, $\tan(4a)=120/119$, and $\tan(4a-b)=1$.
Integrating $1-t^2<1/(1+t^2)<1$ on positive intervals gives
\[
 0<4(1/5-1/(3\cdot5^3))-1/239<4a-b<4/5<1<\pi/2.
\]
Here $\pi/2>1$ follows from $\arctan(1)>1/2$.
Thus $4a-b=\pi/4$ and
\[
 \pi>16(1/5-1/(3\cdot5^3))-4/239
       =\frac{281476}{89625}>\frac{157}{50}>\frac{25}{8}.
\]
This also proves $\phi(0)<2/5$ and $1/(4\pi)<2/25$.
Since $\sqrt{10}<19/6$,
\[
 \CE^2=\frac{19+6\sqrt{10}}{72\pi}
 <\frac{1900}{11304}<\frac{1681}{10000}.
\]
The last difference is $2024/113040000>0$.
\end{proof}

\begin{lemma}[Scalar entry from sharp predecessors]\label{lem:gates}
Put $c=\CE$. If $n\ge47$, $C_{n-1},C_{n-2}\le c$ and $C_n\ge c$, then
\[
 d_n:=n(C_{n-1}-C_n)_+=0,\qquad
 0\le n(C_n-c)<\frac{335909}{1150000}<\frac3{10},
\]
and
\[
 C_n<.4247,\qquad C_{n-1},C_{n-2}<.4253.
\]
If $C_n>c$, the lower bound on the excess is strict.
\end{lemma}
\begin{proof}
Set $s=\sqrt{n/(n-1)}>1$. From \eqref{eq:smooth} and
$s-1=1/((n-1)(s+1))$,
\[
 n(C_n-c)\le\frac n{n-1}
       \left(\frac c{s+1}+\frac s{15}
                         +\frac{2s}{25\sqrt{n-1}}\right).
\]
For all $n\ge47$, $n/(n-1)\le47/46$, $\sqrt{n-1}>6$ and
$s<1011/1000$; the last assertion follows from
$46\cdot1011^2-47\cdot1000^2=17566>0$. Hence
\[
 n(C_n-c)<
 \frac{47}{46}\left(\frac{41}{200}
                  +\frac{1011}{15000}+\frac{1011}{75000}\right)
 =\frac{335909}{1150000}<\frac3{10}.
\]
The margin below $3/10$ is $9091/1150000$. Also
\[
 C_n<\frac{41}{100}+\frac3{470}
     =\frac{1957}{4700}<\frac{4247}{10000},
 \qquad
 C_{n-j}\le c<\frac{41}{100}<\frac{4253}{10000}\quad(j=1,2).
\]
The own-index-cap margin is $3909/470000$. The predecessor sign gives
$d_n=0$. Every comparison includes the endpoint $n=47$.
\end{proof}

\noindent\textbf{Conditional exclusion interface.}
For the next two propositions, an exclusion input includes: attainment
of an oriented positive closed extremum; acquisition of all required
same-law support, probability, moment and resource bounds; and the complete
local exclusion, with every sign, normalization, parameter-domain and
numerical-certificate hypothesis discharged.
Scalar entry alone does not establish this interface.

\begin{proposition}[Weak global assembly]\label{prop:weak}
Assume $C_j\le c=\CE$ for $1\le j\le46$. Suppose the exclusion interface
excludes every selected owner at an index $n\ge47$ satisfying
\[
 0<C_n-c\le.7/n,\quad d_n\le.19,\quad
 C_n<.4247,\quad C_{n-1},C_{n-2}<.4253.
\]
Then $C_n\le c$ for every positive integer $n$.
\end{proposition}
\begin{proof}
If some $C_n>c$, choose the least such $n$. Then $n\ge47$ and both
predecessors are at most $c$. Lemma~\ref{lem:gates} supplies all the
scalar gates, contradicting the stipulated exclusion.
\end{proof}

\begin{proposition}[Separately named equality-inclusive strict tail]
\label{prop:strict}
Assume the same weak prefix through $46$. Suppose the exclusion interface
instead applies also at equality, replacing $0<C_n-c$ above by
$0\le C_n-c$, and excludes those owners as well.
Then $C_n<c$ for every $n\ge47$.
\end{proposition}
\begin{proof}
If the set $\{n\ge47:C_n\ge c\}$ is nonempty, take its least member.
All earlier tail values are $<c$; prefix values are $\le c$.
Both predecessors are therefore at most $c$.
Lemma~\ref{lem:gates} again supplies the scalar gates, including when
$C_n=c$, contradicting the equality-inclusive exclusion.
\end{proof}

The two propositions have different hypotheses and contrary sets.
Proposition~\ref{prop:weak} does not imply strictness: the constant
sequence $C_n=c$ itself satisfies the scalar recurrence.
For Proposition~\ref{prop:strict}, the least $\ge c$ failure is chosen
only within the tail; equality inside the weak prefix is permitted.

\smallskip
\noindent\textbf{Source and verification boundary.}
At the frozen baseline, the recurrence is
\path{paper/analytic-supplement.txt}, physical lines 6479--6522
(label \texttt{cut49:global:smoothing}, original M7458--7506).
The existing tail assembly is \texttt{paper/main.tex}, lines 460--510.
The supplied attainment proof expressly includes equality at lines
107--136; supplement lines 719--755 also permit zero excess in a support
refinement. These observations do not establish equality compatibility
of every intervening acquisition edge. The local comparison and its
binomial input are at supplement lines 7153--7197 and 7748--7768.

The accompanying \texttt{check\_scalar.py} uses only exact integers and
rational arithmetic to check the finite comparisons. Its fresh run
\texttt{runs/least-counterexample-20261027/scalar01/} returned zero;
all thirteen negative controls rejected. That arithmetic PASS proves
neither the analytic steps nor the exclusion interface. This note
identifies a conditional replacement for scalar-entry computations,
not a certificate-deletion list.
\end{document}
